\section{Principal axes, eigenvalues, and rotation}

The immediate problem is the mixed term: changing $i_d$ can alter the
quadratic growth associated with $i_q$, so the two coordinates cannot be read
independently.  We therefore ask a more purposeful question than merely
requesting eigenvalues: is there a direction whose image under $\Hpsm$ changes
only in magnitude, not in direction?  Such a direction satisfies
\begin{equation}
 \Hpsm\vect q=\lambda\vect q.
\end{equation}
Along it, the quadratic form grows as a single scalar coefficient
$\lambda$.  Finding two orthogonal directions of this kind removes the
coupling and reveals the geometric axes.  Without them, one could still use
the double-angle formula below, but the relation between axis direction,
directional voltage sensitivity, and quadratic curvature would remain hidden.

The spectral theorem for real symmetric matrices gives an orthogonal matrix
$\matr Q=[\vect q_1\ \vect q_2]$ and
\begin{equation}
 \Hpsm=\matr Q\matr\Lambda\matr Q\trans,
 \qquad \matr\Lambda=\diag(\lambda_1,\lambda_2),\quad\lambda_i>0.
\end{equation}
With the translated, rotated coordinates
\begin{equation}
 \vect z=\matr Q\trans(\current-\ic),\label{eq:principal-coordinates}
\end{equation}
orthogonality preserves Euclidean lengths and
\begin{equation}
 \lambda_1z_1^2+\lambda_2z_2^2=\Umax^2.\label{eq:principal-ellipse}
\end{equation}
No cross term remains because eigenvectors diagonalize the quadratic map.  A
displacement $t\vect q_j$ changes the left side by $\lambda_jt^2$ alone.  The
eigenvectors are therefore exactly the principal-axis directions, and
\begin{equation}
 a_j=\sqrt{\frac{\Umax^2}{\lambda_j}}=
 \frac{\Umax}{\sqrt{\lambda_j}}.\label{eq:semiaxes}
\end{equation}
A large eigenvalue means steep quadratic growth and hence a short semi-axis.

\paragraph{Normalized thought experiment.}
Suppose the centered equation is
$z_1^2+4z_2^2=1$.  Then $\lambda_1=1$ and $\lambda_2=4$, while
$a_1=1$ and $a_2=1/2$.  Doubling the directional gain squares the
quadratic coefficient and halves the allowed current distance.  This small
example makes the inverse relation between eigenvalue and axis length visible
before machine parameters obscure it.

The \gls{svd} makes the same result especially physical.  Write
\begin{equation}
 \Apsm=\matr U\matr\Sigma\matr V\trans.
\end{equation}
Then
\begin{equation}
 \Hpsm=\Apsm\trans\Apsm
 =\matr V\matr\Sigma^2\matr V\trans.
\end{equation}
Thus the right singular vectors are current-space principal directions and
$\lambda_j=\sigma_j^2$.  A unit current displacement along $\vect v_j$
creates voltage magnitude $\sigma_j$; the allowable distance before reaching
the voltage circle is therefore
\begin{equation}
 a_j=\frac{\Umax}{\sigma_j}.
\end{equation}
The eigenvalue view describes curvature of the pulled-back quadratic form; the
SVD view describes directional gain of the physical current-to-voltage map.
They are not competing explanations but the identity $\Hpsm=\Apsm\trans\Apsm$
seen from two sides.

\begin{figure}[tb]
  \centering
  \begin{tikzpicture}[x=1cm,y=1cm,>=Latex,font=\scriptsize]
    \foreach \x in {0,3.4,6.8,10.2}{
      \draw[gray!35,rounded corners] (\x-1.15,-1.1) rectangle (\x+1.15,1.1);
    }

    \draw[->,gray] (-.9,0)--(.9,0);
    \draw[->,gray] (0,-.9)--(0,.9);
    \draw[current limit] (0,0) circle (.65);
    \node[align=center] at (0,-1.42) {current circle\\isotropic budget};

    \begin{scope}[shift={(3.4,0)},rotate=28]
      \draw[->,gray] (-.9,0)--(.9,0);
      \draw[->,gray] (0,-.9)--(0,.9);
      \draw[current limit] (0,0) circle (.65);
    \end{scope}
    \node[align=center] at (3.4,-1.42) {$\matr V\trans$: rotation\\shape unchanged};

    \draw[->,gray] (5.9,0)--(7.7,0);
    \draw[->,gray] (6.8,-.9)--(6.8,.9);
    \draw[voltage limit] (6.8,0) ellipse [x radius=.9,y radius=.48];
    \node[align=center] at (6.8,-1.42) {$\matr\Sigma$: scaling\\circle becomes ellipse};

    \begin{scope}[shift={(10.2,0)},rotate=-24]
      \draw[->,gray] (-.95,0)--(.95,0);
      \draw[->,gray] (0,-.9)--(0,.9);
      \draw[voltage limit] (0,0) ellipse [x radius=.9,y radius=.48];
    \end{scope}
    \node[align=center] at (10.2,-1.42) {$\matr U$: rotation\\output orientation};

    \draw[->,thick] (1.2,0)--(2.2,0)
      node[midway,above] {$\matr V\trans$};
    \draw[->,thick] (4.6,0)--(5.6,0)
      node[midway,above] {$\matr\Sigma$};
    \draw[->,thick] (8.0,0)--(9.0,0)
      node[midway,above] {$\matr U$};
  \end{tikzpicture}
  \caption{Visual derivation of the SVD. Rotations change orientation but not
  shape or size; only the singular-value scaling changes a circle into an
  ellipse. For the voltage-limit inverse image, the same steps are traversed
  backward, so the allowable current axes scale as $U_{\max}/\sigma_j$.}
  \label{fig:psm-svd-pipeline}
\end{figure}


\subsection{Five viewpoints on the same voltage ellipse}

\begin{center}
\small
\begin{tabularx}{\linewidth}{@{}>{\raggedright\arraybackslash}p{.25\linewidth}X@{}}
\toprule
Viewpoint & What it contributes and what it leaves implicit\\
\midrule
Scalar expansion &
Audits every resistance, speed, inductance, and PM-flux term.  It is best for
sign checking, but the axes are difficult to see.\\
Conic classification &
Uses definiteness, the discriminant, and the translated level to prove that
the boundary is a real ellipse.  It names the shape but says little about the
physical map that created it.\\
Affine inverse image &
Starts with the inverter's voltage circle and pulls it through
$\voltage=\Apsm\current+\vect r$.  It makes the center and boundedness
intuitive, but does not by itself give numerical axis directions.\\
Eigenvalue decomposition &
Diagonalizes $\Hpsm$ and exposes uncoupled current coordinates, curvature, and
semi-axis lengths.  Squaring in the Gram matrix hides the output-side
direction of each voltage response.\\
Singular-value decomposition &
Keeps both current- and voltage-space directions.  It is the clearest gain and
numerical-conditioning interpretation, although it is less convenient for
reading the original scalar coefficients.\\
\bottomrule
\end{tabularx}
\end{center}

All five agree because they start from the identity
\begin{equation}
 \norm{\Apsm\current+\vect r}^2
 =\current\trans\Apsm\trans\Apsm\current
  +2\vect r\trans\Apsm\current+\norm{\vect r}^2.
\end{equation}
Expansion reads the identity in components, conic theory studies its level
set, the inverse-image view studies the map before squaring, the
eigendecomposition studies its Gram matrix, and the SVD factors the original
map.

\begin{figure}[tb]
\centering
\begin{tikzpicture}[scale=1.05]
 \draw[->] (-3.2,0)--(3.4,0) node[right]{$i_d$};
 \draw[->] (0,-2.3)--(0,2.5) node[above]{$i_q$};
 \begin{scope}[shift={(.35,-.15)},rotate=27]
  \draw[constraint curve] (0,0) ellipse (2.35 and 1.05);
  \draw[principal axis] (-2.8,0)--(2.8,0);
  \draw[principal axis] (0,-1.5)--(0,1.5);
  \draw[eigenvector] (0,0)--(1.45,0) node[above]{$\vect q_1$};
  \draw[eigenvector] (0,0)--(0,.8) node[left]{$\vect q_2$};
 \end{scope}
 \node[psm center,label=below right:$\ic$] at (.35,-.15) {};
 \draw[->] (1.2,0) arc[start angle=0,end angle=27,radius=1.2];
 \node at (1.42,.32){$\theta$};
\end{tikzpicture}
\caption{The orthonormal eigenvectors of the quadratic matrix are the ellipse's
principal directions.}
\label{fig:eigenvectors-axes}
\end{figure}

\begin{figure}[tb]
\centering
\begin{tikzpicture}
 \begin{axis}[axis main,width=.78\linewidth,height=5.2cm,
 xlabel={eigenvalue $\lambda$},ylabel={semi-axis $a/U_{\max}$},
 xmin=.2,xmax=5,ymin=0,ymax=2.4,domain=.2:5,samples=100]
  \addplot[very thick,MidnightBlue]{1/sqrt(x)};
  \addplot[only marks,mark=*,BrickRed] coordinates {(.5,1.414) (1,1) (4,.5)};
  \node[annotation,anchor=west] at (axis cs:2.1,1.35)
    {large $\lambda$\\$\Rightarrow$ short axis};
 \end{axis}
\end{tikzpicture}
\caption{Inverse square-root relation between curvature eigenvalue and
semi-axis length.}
\label{fig:eigenvalue-axis}
\end{figure}


The classical angle formula is the same diagonalization written componentwise.
For $Ax^2+Bxy+Cy^2$, substitute
$x=\xi\cos\theta-\eta\sin\theta$ and
$y=\xi\sin\theta+\eta\cos\theta$.  The $\xi\eta$ coefficient becomes
\begin{equation}
 B\cos(2\theta)+(C-A)\sin(2\theta).
\end{equation}
The double angle appears because products of two rotated coordinates contain
$\sin\theta\cos\theta$, $\cos^2\theta$, and $\sin^2\theta$.  Setting the mixed
coefficient to zero gives
\begin{equation}
 \boxed{\tan(2\theta)=\frac{B}{A-C}.}\label{eq:double-angle}
\end{equation}
The two solutions separated by $\pi/2$ are the two orthogonal eigenvector
directions.  Applying \cref{eq:double-angle} to \cref{eq:Hpsm}, where
$B=2H_{12}$, gives for $L_d\neq L_q$
\begin{equation}
 \tan(2\theta)=\frac{2\Rs}{\omegae(\Ld+\Lq)}.\label{eq:psm-angle}
\end{equation}
At a circular limit the angle is undefined because every direction is a
principal direction; an arbitrary numerical angle then carries no physics.

\geointerpretation Translation selects the bowl's minimum, rotation aligns the
coordinate axes with its uncoupled slopes, and the eigenvalues set its
directional steepness.

\physinterpretation Saliency and resistance together create the cross term.
The PM flux appears in $\vect b$ and therefore translates the ellipse, but it
does not enter $\Hpsm$ and hence does not directly change its shape.

\paragraph{Synthesis.}
Translation locates the ellipse, eigenvectors identify independent directions,
eigenvalues measure how rapidly voltage budget is consumed, and singular
values express the same consumption as a physical current-to-voltage gain.
These conclusions require a symmetric positive-definite Gram matrix; a zero
eigenvalue would open an infinite direction, while a negative eigenvalue could
occur only after leaving the Gram-matrix setting.
